In this section, we detail the \textit{dissertation phase} of the proposal, in
which we pursue the two optimizations that the course-phase semantics make
possible: \textit{lazy reasoning} and \textit{parallel reasoning}, which converge
into \system, the deliverable of the proposal.

\subsection{Lazy Reasoning}
\label{sec:lazy}

With a single structure carrying one equivalence relation per assumption
context, lazy reasoning asks how to exploit the dependencies between these
contexts to avoid redundant work: the conditions that guard them are logically
related, so the relation induced by one often determines much of the relation
induced by another. As these relations form a lattice, two operations arise alongside the \textit{fork} that opens a
context: a \textit{meet}, the closure of their union, which merges facts from
different subgoals; and a \textit{join}, their intersection, which keeps the
facts common to every case of a split. The fork alone, a specific form of
laziness that in effect clones a relation on demand, is what versioned e-graphs
already provide; the goal now is to exploit richer dependencies between contexts,
such as those induced by the semantics of the conditions that separate them, e.g.,
so that reasoning under a condition like $a = b \to c = d$ can reuse the relation
already computed under $a = b$. Laziness reifies an inherited equality only when queried, amortizing the cost
over the queries that use it. We can
evaluate the operations in isolation, against eager counterparts, and end to end
inside one of the aforementioned provers, where case splits arise in proofs
by cases, or an SMT solver, where case-splitting paradigms such as
cube-and-conquer~\cite{heule2011cube} are a natural workload. In either setting
we draw on established benchmarks: the suite of the host prover for the former,
and the SMT-LIB benchmark library~\cite{barrett2016smtlib} for the latter,
restricted to the divisions where case splitting dominates the search.

\subsection{Parallel Reasoning}

Parallel reasoning treats the lattice of equivalence relations as a dependency
graph of tasks, each exploring one relation modulo its dependencies, and asks 
how to schedule that graph across many workers. Three difficulties shape the 
problem.
No task is guaranteed to terminate, since saturating a context need not reach a
fixpoint, so execution must be incremental, each task propagating partial
results for downstream tasks to build on rather than running to completion. The
lattice itself may be unbounded, as the conditions worth exploring are
open-ended, so scheduling must interleave advancing the known relations with
discovering new contexts, each discovery triggering a further scheduling round.
And because
the tasks share one underlying structure, a good schedule must keep
synchronization and contention low. Addressing these statically where the
lattice is known, and dynamically where it is not, is the crux of this
direction. We can evaluate a schedule by its speed-up and scalability as the
number of threads grows, relative to sequential exploration and to an ideal
parallelization of the same dependency graph, drawing on the same benchmarks as
in Section~\ref{sec:lazy}.

\subsection{Smart E-Graphs}

\system\ is the deliverable of the proposal: one e-graph library with all four
proposed extensions. A host prover can adopt it in place of its own structure
and enable only the extensions it needs, so that our artifact can reach
the verification ecosystem rather than one tool. Our extensions can interact in
subtle ways, so their composition is non-trivial. For example,
disequalities should be interpreted per version and be preserved by the
meet and join. We will develop \system\ incrementally, validating it against the
benchmarks above in the provers our prior work already targets, and release the
artifact open source.